% Source: https://artofproblemsolving.com/community/c3h2542768_2005_national_team_round_6
% Stable identity: aops-topic-2542768-q01
\emph{Multiples of $3$ are sweet.} Both operations, $x\mapsto2x$ and $x\mapsto x-12$, send multiples of $3$ to multiples of $3$, and $16$ is not a multiple of $3$. So a sequence starting at a multiple of $3$ never contains $16$.

\emph{All other numbers are not sweet.} A sequence starting in $\{1,\dots,50\}$ stays there: $x\le25$ gives $2x\le50$, and $26\le x\le50$ gives $14\le x-12\le38$. The numbers $16\to32\to20\to40\to28\to16$ form a cycle. Following the rule from the other non-multiples of $3$:
\[1\to2\to4\to8\to16,\qquad5\to10\to20,\qquad7\to14\to28,\qquad11\to22\to44\to32,\]
\[13\to26\to14,\qquad17\to34\to22,\qquad19\to38\to26,\qquad23\to46\to34,\qquad25\to50\to38,\]
\[29\to17,\quad31\to19,\quad35\to23,\quad37\to25,\quad41\to29,\quad43\to31,\quad47\to35,\quad49\to37.\]
These chains contain all $34$ non-multiples of $3$ from $1$ to $50$, and each chain leads to a number already shown to reach $16$. So none of them is sweet (including $16$ itself, whose sequence starts with $16$).

Hence the sweet numbers are exactly the $16$ multiples of $3$ from $3$ to $48$: the answer is $\boxed{16}$.
